\documentclass[11pt]{article}

\usepackage[margin=1.5in]{geometry}
\usepackage{lecnotes}
% \usepackage{mpass}
% \lstset{style=verb,language=mpass}
\input{lmacros}

\newcommand{\course}{15-316: Software Foundations of Security \& Privacy}
\newcommand{\lecturer}{Matt Fredrikson}
\newcommand{\lecdate}{April 21, 2026}
\newcommand{\lecnum}{23}
\newcommand{\lectitle}{Mechanisms and Composition}
\newcommand{\courseurl}{https://15316-cmu.github.io/2026/}

\newcommand{\pr}{\mathrm{Pr}}
\newcommand{\flip}{\mathbf{flip}}
\newcommand{\lap}{\mathbf{Lap}}
\newcommand{\sens}{\Delta}

\newaliascnt{example}{theorem}
\newtheorem{example}[example]{Example}
\aliascntresetthe{example}
\providecommand*{\exampleautorefname}{Example}

\begin{document}

\maketitle

\section{Introduction}

In \href{\courseurl/lectures/22-diffprivacy.pdf}{Lecture~22} we introduced
differential privacy as a rigorous property of a randomized
computation. A function $f$ satisfies $\epsilon$-differential privacy if, for
any two neighboring inputs $m_1 \sim_1 m_2$ and any output $s$, we have
$\pr[f(m_1) = s] \leq e^\epsilon \cdot \pr[f(m_2) = s]$.  The parameter
$\epsilon$ controls a tradeoff between the accuracy of the answer and the
indistinguishability of neighboring inputs, and we showed that randomized
response achieves $\mathrm{ln}(3)$-differential privacy for a single secret
bit.

Randomized response is elegant, but it answers only one very specific kind of
question: ``what is the true value of a single yes--no fact?'', and even then
only with noise proportional to the answer itself.  Most of the useful
computations we would like to perform on sensitive datasets---counting the
number of students who earned an A, computing the average response time of an
API, publishing a histogram of income levels by region---do not fit this
template.  They return numeric values from potentially large ranges, rather
than a single bit, and there is no obvious analogue of ``flipping a coin to
return a random answer'' that preserves utility.  Worse, practitioners rarely
have the luxury of running just one query: statistics are typically released
as a batch, and each additional release erodes the privacy guarantee in some
quantifiable way.

Today's lecture builds up the two ingredients that make differential privacy a
practical design discipline rather than a theoretical curiosity.  We begin
with \emph{sensitivity}, the notion that captures ``how much can one
individual's data change the answer to a query?''  We then introduce the
\emph{Laplace mechanism}, a general-purpose recipe that converts any
bounded-sensitivity numeric query into an $\epsilon$-differentially private
approximation by adding calibrated noise.  Finally, we study the
\emph{composition theorems}, which tell us how the privacy budget accumulates
when several differentially private computations are released together.
These three tools---sensitivity, Laplace, and composition---are what lets us
take a realistic data-release scenario and reason directly about its privacy
cost.

\section{Sensitivity}

To use the Laplace mechanism, we first need a way to measure how much an
individual record can affect the output of a numeric query.  The right notion
is the worst-case change in the query's output over all pairs of neighboring
inputs.

\begin{definition}[Global sensitivity]
\label{def:sens}
Let $f : \mathcal{D} \to \mathbb{R}$ be a real-valued query on memories.  The
\emph{global sensitivity} of $f$ is
\[
  \sens f = \max_{m_1 \sim_1 m_2} |f(m_1) - f(m_2)|.
\]
\end{definition}

Sensitivity is a property of the query, not of any particular dataset; it
quantifies the worst damage that \emph{any} single individual's data can do
to the answer.  Computing it requires thinking carefully about what it means
for two memories to ``differ in exactly one index'', and what the extremal
values of the differing index could be.  We will assume throughout this
lecture that individual values are bounded in some known range, say $[0, U]$,
without which global sensitivity would be infinite for most queries of
interest.

\begin{example}[Counting query]
\label{ex:sens-count}
Let $f_{\geq t}(m) = |\{i : m[i] \geq t\}|$ count the number of cells whose
value exceeds threshold $t$.  Two neighboring memories $m_1$ and $m_2$ differ
in one index; that index either satisfies the threshold in both, neither, or
exactly one.  In the first two cases $f_{\geq t}(m_1) = f_{\geq t}(m_2)$, and
in the last case they differ by exactly one.  Thus
$\sens f_{\geq t} = 1$, regardless of the threshold or the size of the
database.
\end{example}

\begin{example}[Sum and mean]
\label{ex:sens-sum}
Suppose each $m[i] \in [0, U]$.  For the sum
$f_\Sigma(m) = \sum_i m[i]$, neighboring memories can differ by as much as
$U$ at the one differing index, so $\sens f_\Sigma = U$.  For the mean
$f_\mu(m) = (\sum_i m[i])/N$ over $N$ records, $\sens f_\mu = U/N$:
the change in one record is averaged over $N$.
\end{example}

\begin{example}[Max]
\label{ex:sens-max}
Consider the maximum $f_{\max}(m) = \max_i m[i]$, with each $m[i] \in [0, U]$.
Take $m_1 = (U, 0, 0, \ldots, 0)$ and $m_2 = (0, 0, 0, \ldots, 0)$; they are
neighbors, and $f_{\max}(m_1) - f_{\max}(m_2) = U$.  Hence
$\sens f_{\max} = U$.  The max is an intuitively ``leaky'' query, and its
global sensitivity reflects that: a single record can move the answer across
the entire range.
\end{example}

Small-sensitivity queries (counts, bounded averages) can be released with
relatively little noise, while high-sensitivity queries like the max require
more noise to be made private, at the cost of utility.  This is a recurring
theme in differential privacy: the structure of the query itself dictates how
much randomness you must inject.

\section{The Laplace Mechanism}

The Laplace distribution provides a principled way to add noise.  Its
probability density function is
\[
  \lap(b)(x) = \frac{1}{2b}\exp\!\left(-\frac{|x|}{b}\right),
\]
where $b > 0$ is the \emph{scale} parameter.  Larger $b$ means heavier tails
and more variance: the variance of $\lap(b)$ is $2b^2$, and its standard
deviation is $b\sqrt 2$.

The Laplace mechanism adds Laplace noise calibrated to the sensitivity of the
query.

\begin{definition}[Laplace mechanism]
\label{def:laplace}
Let $f : \mathcal{D} \to \mathbb{R}$ have global sensitivity $\sens f$, and
let $\epsilon > 0$.  The \emph{Laplace mechanism} is the randomized
function
\[
  M_\mathrm{Lap}(m) = f(m) + Y, \qquad Y \sim \lap(\sens f / \epsilon).
\]
\end{definition}

The scale is chosen so that a change of $\sens f$ in the query's value is
``absorbed'' by the noise to within the privacy budget, as the following
theorem makes precise.

\begin{theorem}[Privacy of the Laplace mechanism]
\label{thm:laplace}
The Laplace mechanism $M_\mathrm{Lap}$ satisfies $\epsilon$-differential
privacy.
\end{theorem}

\begin{proof}
We need to show that for all $m_1 \sim_1 m_2$ and all $s \in \mathbb{R}$,
$\pr[M_\mathrm{Lap}(m_1) = s] \leq e^\epsilon \cdot \pr[M_\mathrm{Lap}(m_2) = s]$.
Since the output is continuous we work with density ratios, but the
argument is identical.  Let $b = \sens f / \epsilon$.  The density of
$M_\mathrm{Lap}(m)$ at $s$ is $\lap(b)(s - f(m))$, so
\[
  \frac{\pr[M_\mathrm{Lap}(m_1) = s]}{\pr[M_\mathrm{Lap}(m_2) = s]}
  = \frac{\exp(-|s - f(m_1)|/b)}{\exp(-|s - f(m_2)|/b)}
  = \exp\!\left(\frac{|s - f(m_2)| - |s - f(m_1)|}{b}\right).
\]
By the reverse triangle inequality,
$|s - f(m_2)| - |s - f(m_1)| \leq |f(m_1) - f(m_2)| \leq \sens f$, so
\[
  \frac{\pr[M_\mathrm{Lap}(m_1) = s]}{\pr[M_\mathrm{Lap}(m_2) = s]}
  \leq \exp(\sens f / b) = \exp(\epsilon) = e^\epsilon. \qedhere
\]
\end{proof}

The theorem reduces a difficult question (``is this computation private?'')
to a much more tractable one (``what is the sensitivity of my query?'').
Once you know $\sens f$, adding $\lap(\sens f / \epsilon)$ noise is a
mechanical step.  Observe that \emph{smaller} privacy budgets $\epsilon$
demand \emph{larger} noise scales: to cut the budget in half, you double the
noise.  Doubling a query's sensitivity has the same effect.

\begin{example}[Private mean]
\label{ex:private-mean}
Returning to \autoref{ex:sens-sum} with $U = 100$ and $N = 50$ students, the
Laplace mechanism can release the class mean with $\epsilon = 1$ by adding
noise drawn from $\lap(100/(50 \cdot 1)) = \lap(2)$.  The standard deviation
of the noise is $2\sqrt 2 \approx 2.83$, so the private mean is typically
within a few points of the true average.  Contrast this with releasing the
\emph{maximum} grade with $\epsilon = 1$: by \autoref{ex:sens-max} the scale
becomes $\lap(100)$, with standard deviation roughly $141$---the noise is
larger than the entire range of valid answers, rendering the release almost
useless.  This is why practitioners prefer aggregate queries (sums, counts,
means) to order statistics when possible, and reach for more specialized
mechanisms, such as the \emph{exponential mechanism}~\citep{Dwork14fttcs},
when an order statistic or categorical answer is unavoidable.
\end{example}

\section{Composition}

A single private query is rarely enough.  Research studies release dozens of
statistics; dashboards update continuously; machine-learning pipelines make
repeated passes over the same training set.  Each of these releases costs
some privacy, and we need a principled way to track the total.  Three
composition theorems cover the cases that arise most often.

\subsection{Sequential Composition}

The most general fact is that privacy budgets add under independent queries.

\begin{theorem}[Sequential composition]
\label{thm:seq-comp}
Let $f_1$ and $f_2$ be randomized mechanisms satisfying $\epsilon_1$- and
$\epsilon_2$-differential privacy respectively, and suppose they use
independent randomness.  Then the combined mechanism
$f_{12}(m) = (f_1(m), f_2(m))$ satisfies $(\epsilon_1 + \epsilon_2)$-
differential privacy.
\end{theorem}

\begin{proof}
Fix neighboring inputs $m_1 \sim_1 m_2$ and an output pair $(s_1, s_2)$.  By
independence of $f_1$ and $f_2$'s randomness,
\[
  \pr[f_{12}(m_1) = (s_1, s_2)] = \pr[f_1(m_1) = s_1] \cdot \pr[f_2(m_1) = s_2].
\]
Applying the $\epsilon_i$-DP guarantee of each factor,
\[
\begin{array}{rl}
  \pr[f_{12}(m_1) = (s_1, s_2)]
  &\leq\ e^{\epsilon_1}\pr[f_1(m_2) = s_1]\cdot e^{\epsilon_2}\pr[f_2(m_2) = s_2] \\[0.5ex]
  &=\ e^{\epsilon_1+\epsilon_2}\pr[f_{12}(m_2) = (s_1, s_2)]. \qedhere
\end{array}
\]
\end{proof}

Sequential composition is the reason $\epsilon$ is called a ``budget'': each
query spends a portion of it, and once the total exceeds what the data
custodian is willing to grant, no further releases are allowed.  It also
holds \emph{adaptively}: even if the choice of $f_2$ depends on the output of
$f_1$, the total budget is still $\epsilon_1 + \epsilon_2$.  Adaptive
composition is what makes differential privacy compose cleanly with
interactive analyses, where the next question depends on the previous
answer.

\subsection{Parallel Composition}

When queries operate on \emph{disjoint} partitions of the data, we can do
better than simply summing budgets.

\begin{theorem}[Parallel composition]
\label{thm:par-comp}
Let $\{D_1, D_2, \ldots, D_k\}$ be a partition of the index set of $m$ into
disjoint blocks, and let $f_i$ be an $\epsilon_i$-DP mechanism applied to
the projection of $m$ onto $D_i$.  Then the joint release
$(f_1(m|_{D_1}), \ldots, f_k(m|_{D_k}))$ satisfies
$\max_i \epsilon_i$-differential privacy.
\end{theorem}

\begin{proof}
Take neighbors $m_1 \sim_1 m_2$; they differ in exactly one index $j$, which
lies in exactly one block $D_{i^*}$.  For every other block $D_i$ with
$i \neq i^*$, the projections are identical, so $f_i$'s output distribution
is identical on $m_1$ and $m_2$.  For the one block $D_{i^*}$ that contains
$j$, the mechanism $f_{i^*}$ is $\epsilon_{i^*}$-DP.  Combining these two
observations with independence gives the ratio bound $e^{\epsilon_{i^*}}$,
and since $i^*$ could be any index we conclude $\max_i \epsilon_i$.
\end{proof}

Parallel composition is why a histogram over disjoint buckets is not much
more expensive than a single count: an individual can fall into at most one
bucket, so changing their record affects at most one count.  Intuitively, a
partitioned release lets each block ``share'' the privacy budget rather than
additively paying for each one.  The key word is \emph{disjoint}: if the
blocks overlap, parallel composition does not apply, and we must fall back
to sequential composition.

\subsection{Post-Processing}

The third property is perhaps the most reassuring: once a release is
differentially private, no amount of further computation on the released
value can make it less so.

\begin{theorem}[Post-processing]
\label{thm:post-proc}
Let $f$ be an $\epsilon$-DP mechanism and $g$ an arbitrary (possibly
randomized) function that depends only on the output of $f$.  Then
$g \circ f$ is $\epsilon$-DP.
\end{theorem}

\begin{proof}
For any neighbors $m_1 \sim_1 m_2$ and any output $s'$,
\[
\begin{array}{rl}
  \pr[g(f(m_1)) = s']
  &=\ \sum_s \pr[g(s) = s']\,\pr[f(m_1) = s] \\[0.5ex]
  &\leq\ \sum_s \pr[g(s) = s']\, e^\epsilon\, \pr[f(m_2) = s] \\[0.5ex]
  &=\ e^\epsilon \pr[g(f(m_2)) = s']. \qedhere
\end{array}
\]
\end{proof}

Post-processing is essential in practice: we often want to take a noisy
count, round it to an integer, clamp it to a valid range, or feed it into
downstream analytics.  None of these operations can harm the privacy
guarantee, as long as they do not peek back at the raw data.

\subsection{A Worked Example}

Let us return to a realistic scenario.  Suppose the course staff wants to
release the following three statistics about the class of $N$ students, and
must decide how to split a total privacy budget of $\epsilon = 2$:
\begin{enumerate}
\item[(a)] The mean grade.
\item[(b)] The count of students who scored an $A$ (grade in $[90, 101)$).
\item[(c)] The count of students who scored a $C$ or below (grade in
  $[0, 70)$).
\end{enumerate}

Each of (a), (b), (c) can be implemented as a Laplace mechanism with its own
budget.  By \autoref{ex:sens-count} and \autoref{ex:sens-sum}, we have
$\sens f_\mu = 100/N$ and $\sens f_{\geq 90} = \sens f_{<70} = 1$.

Observe that (b) and (c) operate on \emph{disjoint} ranges: a student's
grade is in $[90, 101)$ or in $[0, 70)$ or in $[70, 90)$, and these sets do
not overlap.  Therefore we can apply \autoref{thm:par-comp} to (b) and (c)
together, spending only $\max(\epsilon_b, \epsilon_c)$ on the pair rather
than their sum.  The mean in (a) uses every student's grade, so it overlaps
both (b) and (c); it must be composed sequentially with the pair.

If we set $\epsilon_a = 1$, $\epsilon_b = 1$, $\epsilon_c = 1$, the overall
budget is $\epsilon_a + \max(\epsilon_b, \epsilon_c) = 1 + 1 = 2$, exactly
what we allotted.  The mechanism for (a) adds $\lap(100/N)$ noise; the
mechanisms for (b) and (c) each add $\lap(1)$ noise.  If the staff
additionally post-processes the noisy counts by rounding and clamping to
$[0, N]$, \autoref{thm:post-proc} ensures the overall budget does not
change.

Had we ignored parallel composition and simply added all three budgets, we
would have computed $\epsilon_a + \epsilon_b + \epsilon_c = 3$---a $50\%$
overestimate of the real cost.  Recognizing which queries partition the data
and which overlap is therefore a load-bearing step in any practical privacy
analysis.

\section{Conclusion}

The pieces assembled in this lecture---sensitivity, the Laplace mechanism,
and the three composition theorems---form the standard toolkit for
designing differentially private computations in practice.  The broader
literature has developed many refinements: advanced composition theorems
with tighter bounds for large numbers of queries, the exponential mechanism
for selecting among categorical options with bounded utility loss, and
R\'enyi differential privacy for reasoning about Gaussian noise used in
private machine learning~\citep{Dwork14fttcs}.  Language-based approaches go
further still, building type systems that track $\epsilon$ automatically as
a program is written~\citep{Reed10icfp,Gaboardi13popl}, so that privacy
accounting becomes a compile-time check rather than a manual
bookkeeping exercise.  What all of these techniques share is the same
structure we have seen today: identify what one individual can influence,
calibrate noise to hide that influence, and account carefully for the cost
when multiple computations are released together.

\section{Practice Problems}

\begin{task}\rm
  Compute the global sensitivity of each of the following queries, assuming
  each $X(i) \in [0, 100]$ and there are $N$ students.  Briefly justify your
  answer, in particular by exhibiting a pair of neighbors that witnesses the
  bound.
  \begin{enumerate}
  \item The number of students whose grade is a perfect $100$.
  \item The sum of all grades.
  \item The difference between the mean grade and the median grade.
  \item The number of students whose grade lies strictly within one standard
    deviation of the mean grade.
  \end{enumerate}
\end{task}

\begin{task}\rm
  The TAs want to release four Laplace-noised statistics about 15-316,
  each calibrated for $\epsilon = 0.5$:
  \begin{enumerate}
  \item[(a)] The mean grade over all students in the course.
  \item[(b)] The count of freshmen whose final grade is at least $80$.
  \item[(c)] The count of sophomores whose final grade is at least $80$.
  \item[(d)] The count of students (any class) who submitted Lab 3 on time.
  \end{enumerate}
  \begin{enumerate}
  \item Identify which pairs or subsets of (a)--(d) can be composed in
    parallel, and which must be composed sequentially.  Justify your
    reasoning in terms of which partitions of the class each query consumes.
  \item Compute the tightest total privacy budget $\epsilon$ that follows
    from \autoref{thm:seq-comp} and \autoref{thm:par-comp}.
  \item Suppose the TAs then post-process the noisy count in (d) by dividing
    it by the total enrollment to report a ``fraction of students on time''
    instead of a count.  Does this post-processing change the privacy
    budget?  State which composition theorem justifies your answer.
  \end{enumerate}
\end{task}

\bibliographystyle{plainnat}
\bibliography{bibliography}

\end{document}
